\section{Interpretation and recommendations}
\label{sec:discussion}

\subsection{What the results mean}

\paragraph{A simple physical rule, a strong personal component.} Young adults clear an obstacle by lifting each foot by about the
obstacle height (about 10\,mm per cm) while keeping their safety margin, mainly through leading-thigh flexion and trailing heel-up.
On top of this rule every person has a large personal component: a walking signature that identifies them from a single crossing,
and a personal obstacle margin that varies between people by two thirds of the step between two heights. This explains the
central result of the machine learning: a model that has never seen the person confuses neighbouring heights (generic models
56--73\,\%), while scaling the features with the person's own crossings, without any label, removes the personal component and
gives 90--92\,\% for a new person.

\paragraph{What the wearables can and cannot recognise.} Obstacle \emph{height} is recognisable, as a class and as a number
(about 2.4\,cm error), with two foot IMUs, from the leading limb early enough to cue the trailing limb, and without any camera.
Obstacle \emph{width} (5 vs 15\,cm) is not, and the reason is behavioural: people barely change how they cross, so neither the
camera foot placement nor the side-to-side motion separates the widths. The \emph{toe clearance} of a crossing depends mostly on
how high the toe itself goes; the foot sensor tracks this only partly, so clearance estimation needs a better toe-height estimate
(Section~\ref{sec:res-clearance}).

\paragraph{Robustness of the evaluation.} Every accuracy is for an unseen participant, every fitted step stays inside the training
participants, the main settings were also chosen by nested selection, permutation tests reject chance ($p \leq 0.003$), and an
independent Python implementation of the main models agrees with MATLAB (logistic regression: mean absolute difference 0.0 points
over 76 results, maximum 0.4; random forest: 0.9 points over 39 results, maximum 2.2, because the two libraries grow different
trees). The person scaling uses each person's own crossings
without their labels, which is exactly the information a device collects while it is worn.

\subsection{Recommendations for a real-time haptic cue}
\label{sec:haptic-reco}

Table~\ref{tab:design} answers the design questions with the evidence behind each answer.

\begin{table}[htbp]
\centering
\caption{Design of a haptic cue for the trailing limb: answers and evidence (deadline = trailing toe lift-off).}
\label{tab:design}
\footnotesize
\begin{tabular}{@{}p{2.6cm} p{7.4cm} p{5.6cm}@{}}
\toprule
\textbf{Question} & \textbf{Answer} & \textbf{Evidence} \\
\midrule
When to cue? & At the latest at the leading-foot apex, about 0.45\,s before the trailing toe leaves the floor. Detect the crossing earlier
(10--30\,\% of the leading swing, 0.57--0.71\,s before) and decide the height by the apex. Do not wait for the leading landing
(0.2\,s left; in 22\,\% of crossings the trailing toe has already lifted). &
Table~\ref{tab:timeline}, \ref{tab:cueing}, \ref{tab:detect}; height about 80\,\% known at the apex (Table~\ref{tab:timing}). \\
\addlinespace
Which limb? & The \textbf{leading} limb gives the cue information. The trailing limb shows no height information before it leaves the
floor; its strong signals (heel-up shank, foot lift) come 0.25--0.3\,s after lift-off and can only confirm the crossing or give feedback
afterwards. & Table~\ref{tab:realtime} (trailing 41--51\,\% before lift-off), Table~\ref{tab:grade}. \\
\addlinespace
Which segment? & A \textbf{foot} IMU on each foot (foot height and trajectory; the leading leg varies). If only angle sensors are used:
the \textbf{leading thigh}. The leading shank carries no height information. & Tables~\ref{tab:sensors}, \ref{tab:segments},
\ref{tab:grade}, \ref{tab:map}. \\
\addlinespace
Detection threshold? & Leading thigh flexion more than about 6\textdegree{} above the person's normal walking, or leading foot more than
about 19\,mm above it, early in the leading swing. & Table~\ref{tab:detect}: 97--99\,\% sensitivity, 98\,\% specificity. \\
\addlinespace
Height thresholds? & \textbf{Personal}: two boundaries around the person's typical crossing, e.g.\ leading-foot apex height $-36$ and
$+39$\,mm, or leading thigh at 72\,\% of its swing $-3$\textdegree{} and $+5$\textdegree{}. Population thresholds referenced only to
normal walking work poorly. & Table~\ref{tab:grade} (73\,\% / 64\,\% personal vs 45\,\% / 52\,\% population), Table~\ref{tab:margin}. \\
\addlinespace
Calibration? & A short normal walk plus the person's own crossings (no labels): 10--15 crossings give 81--85\,\%, a session 92\,\%;
or one labelled crossing (80\,\%). Keep updating the scaling during use. & Tables~\ref{tab:calibration}, \ref{tab:mixed}. \\
\addlinespace
Reliability? & Cue only when the model is confident (probability at least 0.6): 79--82\,\% of the given cues are right; the stated
probabilities understate the model's real accuracy. & Tables~\ref{tab:cue}, \ref{tab:confidence}, \ref{tab:reliability}. \\
\addlinespace
How much lift? & Per height H1/H2/H3, relative to normal walking: trailing foot +237/+352/+436\,mm (10.9\,mm per cm), trailing shank
heel-up $-31$/$-46$/$-55$\textdegree{} ($-1.3$\textdegree{} per cm); crossing-to-crossing SD 36\,mm and 5.6\textdegree{}. &
Table~\ref{tab:dose}. \\
\bottomrule
\end{tabular}
\end{table}

\paragraph{Reaction time.} The time available after a cue at the leading-foot apex is about 0.45\,s (median; 0.35\,s for a quarter of
the crossings). A simple reaction to a vibration typically takes in the order of 0.2\,s in young adults and longer in older adults;
this was not measured here, and must be measured in the target users before the cue moment is fixed. If it is too long, the earlier
crossing detection (0.57--0.71\,s before the deadline) can trigger a generic ``obstacle'' warning, followed by the height-specific
cue at the apex.

\paragraph{A ``lift more'' feedback.} If the cue should also tell the trailing limb how much to lift, a natural target is the
person's own mean lift for the predicted height, and an alarm threshold about one to two crossing-to-crossing SDs below it
(Table~\ref{tab:dose}). This is a design proposal that has not been tested.


\subsection{Limitations}
\label{sec:limitations}

\begin{itemize}
  \item \textbf{Ten young adults.} All results are feasibility results; they need confirmation on new participants and on the
  target users (older adults), whose crossing strategies and reaction times differ.
  \item \textbf{Blocked obstacle order.} Real sessions mix obstacles; the mixed order was simulated by shuffling the recorded
  crossings, which assumes that a crossing does not depend on the previous obstacle.
  \item \textbf{Camera-defined roles and strides.} Most analyses take the leading/trailing roles and crossing strides from the camera;
  the IMU recovers them (100\,\% role agreement on the detected crossings, 95\,\% fully paired in the camera-free pipeline), and the
  camera-free pipeline gives the end-to-end figure.
  \item \textbf{Feature design after the first results.} The clearance-shape features were designed after the first models, guided
  by biomechanics; a small optimism may remain in their generic-model gain.
  \item \textbf{Trailing toe-off event.} The step-4 toe-off of the trailing limb comes about 0.23\,s after the toe leaves the floor; the
  trailing swing features are consistently defined by that event, and all timing statements use the camera lift-off instead.
  \item \textbf{Offline simulation of real time.} The real-time cue is simulated on recorded data; the reaction time to a vibration and
  the effect of the cue on the crossing were not measured.
\end{itemize}

\subsection{Next steps}

\begin{enumerate}
  \item Record more participants, including older adults, and run the obstacles in a mixed order.
  \item Implement the cue on the device (two foot IMUs, leading-limb trigger) and measure the reaction time and the effect of the cue
  on the trailing-foot clearance.
  \item Measure the foot-sensor position on the shoe so the toe height can be computed instead of fitted, to estimate clearance.
  \item Test obstacles of other heights and depths (e.g.\ kerbs and steps) to confirm the generalisation to unseen heights.
\end{enumerate}
